\documentclass[aspectratio=169]{beamer}
\usetheme{CambridgeUS}
\usecolortheme{beaver}

\newcommand\hmmax{0}
\newcommand\bmmax{0}

\usepackage{stmaryrd}

\usepackage{amsmath,amssymb,enumerate,amsthm}
\usepackage{bm}
\usepackage{graphicx}

\usepackage{algorithm, algpseudocode}
\usepackage[strict]{siunitx}
\usepackage{hyperref}
\usepackage{mathrsfs}

\usepackage{bm}
\usepackage{empheq}
%\usepackage{cleveref}
\usepackage{xcolor}
\usepackage{textcomp}
\usepackage{siunitx}
\usepackage{booktabs}
\usepackage{caption}
\usepackage{subcaption}
\usepackage{tikz}
\usetikzlibrary{decorations.pathreplacing,calligraphy,graphs}
\usepackage{pgfplots}
\pgfplotsset{compat=1.17}
\usepackage{lmodern}
\usepackage{array}
\usepackage{mathtools}
\usepackage{multicol}
\newcolumntype{C}[1]{>{\centering\arraybackslash}m{#1}}
\NewDocumentCommand{\R}{}{\mathbb{R}}
\NewDocumentCommand{\N}{}{\mathbb{N}}
%\mathscr is rounder than \mathcal.
\NewDocumentCommand{\powerset}{m}{\mathscr{P}(#1)}
%Powerset without zero.
\NewDocumentCommand{\powersetz}{m}{\mathscr{P}^*(#1)}
%https://tex.stackexchange.com/a/45732, works within both \set and \set*, same spacing than \mid (https://tex.stackexchange.com/a/52905).
\NewDocumentCommand{\suchthat}{}{\;\ifnum\currentgrouptype=16 \middle\fi|\;}
%Integer interval.
\NewDocumentCommand{\intvl}{m}{\llbracket#1\rrbracket}
%Allows for \abs and \abs*, which resizes the delimiters.
\DeclarePairedDelimiter\abs{\lvert}{\rvert}
\DeclarePairedDelimiter\card{\lvert}{\rvert}
\DeclarePairedDelimiter\floor{\lfloor}{\rfloor}
\DeclarePairedDelimiter\ceil{\lceil}{\rceil}
%Perhaps should use U+2016 ‖ DOUBLE VERTICAL LINE here?
\DeclarePairedDelimiter\norm{\lVert}{\rVert}
%From mathtools. Better than using the package braket because braket introduces possibly undesirable space. Then: \begin{equation}\set*{x \in \R^2 \suchthat \norm{x}<5}\end{equation}.
\DeclarePairedDelimiter\set{\{}{\}}
\DeclareMathOperator*{\argmax}{arg\,max}
\DeclareMathOperator*{\argmin}{arg\,min}

%UTR #25: Unicode support for mathematics recommend to use the straight form of phi (by default, given by \phi) rather than the curly one (by default, given by \varphi), and thus use \phi for the mathematical symbol and not \varphi. I however prefer the curly form because the straight form is too easy to mix up with the symbol for empty set.
\let\phi\varphi

%The amssymb solution.
%\NewDocumentCommand{\restr}{mm}{{#1}_{\restriction #2}}
%Another acceptable solution.
%\NewDocumentCommand{\restr}{mm}{{#1|}_{#2}}
%https://tex.stackexchange.com/a/278631; drawback being that sometimes the text collides with the line below.
\NewDocumentCommand\restr{mm}{#1\raisebox{-.5ex}{$|$}_{#2}}


%Decision Theory (MCDA and SC)
\NewDocumentCommand{\allalts}{}{\mathcal{A}}
\NewDocumentCommand{\allcrits}{}{\mathscr{C}}
\NewDocumentCommand{\alts}{}{A}
\NewDocumentCommand{\dm}{}{i}
\NewDocumentCommand{\allF}{}{\mathscr{F}}
\NewDocumentCommand{\allvoters}{}{\mathcal{N}}
\NewDocumentCommand{\voters}{}{N}
\NewDocumentCommand{\prof}{}{P}
\NewDocumentCommand{\linors}{}{\mathcal{L}(\allalts)}
\NewDocumentCommand{\alllosses}{}{\intvl{0, m-1}^N}
%Thanks to https://tex.stackexchange.com/q/154549
	%\makeatletter
	%\def\@myRgood@#1#2{\mathrel{R^X_{#2}}}
	%\def\myRgood{\@ifnextchar_{\@myRgood@}{\mathrel{R^X}}}
	%\makeatother
\NewDocumentCommand{\pref}{}{\succ}
\NewDocumentCommand{\prefi}{O{i}}{\succ_{#1}}
\NewDocumentCommand{\PO}{}{\mathit{PO}}

%Deliberated Judgment
\NewDocumentCommand{\allargs}{}{S^*}
\NewDocumentCommand{\args}{}{S}
\NewDocumentCommand{\ar}{}{s}
\NewDocumentCommand{\allprops}{}{T}
\NewDocumentCommand{\prop}{}{t}
\NewDocumentCommand{\ileadsto}{}{⇝}
\NewDocumentCommand{\ibeatse}{}{⊳_\exists}
\NewDocumentCommand{\nibeatse}{}{⋫_\exists}
\NewDocumentCommand{\ibeatsst}{}{⊳_\forall}
\NewDocumentCommand{\nibeatsst}{}{⋫_\forall}
\NewDocumentCommand{\mleadsto}{O{\eta}}{⇝_{#1}}
\NewDocumentCommand{\mbeats}{O{\eta}}{⊳_{#1}}
\NewDocumentCommand{\ibeatseinv}{}{⊳_\exists^{-1}}

%Logic
\NewDocumentCommand{\ltru}{}{\texttt{T}}
\NewDocumentCommand{\lfal}{}{\texttt{F}}
\newtheorem{remark}{Remark}

\newcommand{\profile}{\bm{v}}%(complete) profile
\newcommand{\pprofile}{{\bm{p}}}%partial profile
\newcommand{\w}{\bm{w}}
\newcommand{\W}{\mathcal{W}}
\newcommand{\Co}{\mathcal{C}}
\newcommand{\pw}{W}%our knowledge about the weights
\newcommand{\strat}[1]{\emph{#1}}
\newcommand{\ppref}{\succ^\text{p}}%partial pref
\newcommand{\pprefeq}{\succeq^\text{p}}%partial pref
\DeclareMathOperator{\Regret}{Regret}
\DeclareMathOperator{\SCORE}{Score}
\DeclareMathOperator{\PMR}{PMR}
\DeclareMathOperator{\MR}{MR}
\DeclareMathOperator{\MMR}{MMR}

\newcommand*{\icimg}[1]{%
	\raisebox{-.3\baselineskip}{%
		\includegraphics[
		height=\baselineskip,
		width=\baselineskip,
		keepaspectratio,
		]{#1}%
	}%
}

\newcommand*{\icarr}[1]{%
	\raisebox{-0.4\baselineskip}{%
		\includegraphics[
		height=2.5\baselineskip,
		width=3\baselineskip,
		keepaspectratio,
		]{#1}%
	}%
}

%%Title page
\makeatletter
\defbeamertemplate*{title page}{mytitle}[1][]
{
	\vbox{}
	\vfill
	\begin{centering}

%{\usebeamercolor[fg]{titlegraphic}\inserttitlegraphic\par}
		\begin{beamercolorbox}{titlegraphic}
				\usebeamerfont{titlegraphic}\inserttitlegraphic
		\end{beamercolorbox}%
			\vskip1em\par	
		\begin{beamercolorbox}[rounded=true, center, shadow=true, sep=8pt,#1]{title}
			\usebeamerfont{title}\inserttitle\par%
			\ifx\insertsubtitle\@empty%
			\else%
			\vskip0.5em%
			{\usebeamerfont{subtitle}\usebeamercolor[fg]{subtitle}\insertsubtitle\par}%
			\fi%     
		\end{beamercolorbox}%
		\vskip1em\par
		\begin{beamercolorbox}[sep=8pt,center,#1]{author}
			\usebeamerfont{author}\insertauthor
		\end{beamercolorbox}
		\begin{beamercolorbox}[sep=8pt,center,#1]{institute}
			\usebeamerfont{institute}\insertinstitute
		\end{beamercolorbox}
		\begin{beamercolorbox}[sep=8pt,center,#1]{date}
			\usebeamerfont{date}\insertdate
		\end{beamercolorbox}\vskip0.5em
		\begin{beamercolorbox}[sep=8pt,center,#1]{logo}
			\insertlogo
		\end{beamercolorbox}%
	\end{centering}
	
}
\setbeamertemplate{title page}[mytitle]
\makeatother

%Footline
\makeatletter
\setbeamercolor{progress bar progress}{use=progress bar,bg=progress bar.fg}

\def\progressbar@progressbar{} % the progress bar
\newcount\progressbar@tmpcounta% auxiliary counter
\newcount\progressbar@tmpcountb% auxiliary counter
\newdimen\progressbar@pbht %progressbar height
\newdimen\progressbar@pbwd %progressbar width
\newdimen\progressbar@tmpdim % auxiliary dimension
\progressbar@pbwd=\paperwidth
\progressbar@pbht=0.5ex
% the progress bar

\defbeamertemplate{footline}{progress bar}{
	\progressbar@tmpcounta= \insertframenumber % max = ?
	\progressbar@tmpcountb=\inserttotalframenumber      
	\progressbar@tmpdim=.5\progressbar@pbwd
	\multiply\progressbar@tmpdim by \progressbar@tmpcounta
	\divide\progressbar@tmpdim by \progressbar@tmpcountb
	\progressbar@tmpdim=2\progressbar@tmpdim

	\leavevmode%
	\begin{beamercolorbox}[wd=\paperwidth,ht=2ex,dp=1ex,right]{bar}
		\insertframenumber{} / \inserttotalframenumber\hspace*{2ex} 
	\end{beamercolorbox}
	\newline
	\begin{beamercolorbox}[wd=\paperwidth,ht=1ex,dp=1ex]{progress bar}
		\begin{beamercolorbox}[wd=\progressbar@tmpdim,ht=1ex,dp=1ex]{progress bar progress}
		\end{beamercolorbox}%
	\end{beamercolorbox}%
	
	
}
\setbeamertemplate{footline}[progress bar]
\setbeamercolor{progress bar}{fg=red!65!black,bg=white!85!black}
\makeatother

%Headline empty
\setbeamertemplate{headline}{}



\titlegraphic{\includegraphics[width=50mm]{logo_dauphine} \hspace*{9cm} \includegraphics[width=7mm]{cnrs}}

\title[]{Simultaneous Elicitation of Scoring Rule and Agent Preferences for Robust Winner Determination}
%
%\subtitle{Proposal: ``Elicitation and Explanation for Voting Rules''}
\institute[]{\textsuperscript{1} Université Paris-Dauphine, Université PSL, CNRS, LAMSADE \\ \textsuperscript{2} LIP6, UMR 7606, CNRS and Sorbonne Universit\'e}
\author{\texorpdfstring{\textbf{B. Napolitano}\textsuperscript{1}, O. Cailloux\textsuperscript{1} and P. Viappiani\textsuperscript{2}}{}}
\date{\small{ADT 2021 - 04 November 2021} \\ \includegraphics[width=35mm]{LOGO_LAMSADE} }

\usepackage{tikz}
\usepackage{amsmath}
\usepackage{graphicx}

\definecolor{darkred}{rgb}{0.8,0,0}

\begin{document}

\beamertemplatenavigationsymbolsempty

\begin{frame}[plain]
\maketitle
\end{frame}

\addtocounter{framenumber}{-1}

\begin{frame}
	\frametitle{Introducing the problem}
		\vspace{-2em}
		\small{Agents} $=\{$ \includegraphics[scale=0.05]{assets/v1.png}, \includegraphics[scale=0.05]{assets/v2.png}, \includegraphics[scale=0.05]{assets/v3.png} $\}$, $\ $ \small{Alternatives} $=\{$ \includegraphics[scale=0.025]{assets/d1.png}, \includegraphics[scale=0.025]{assets/d2.png}, \includegraphics[scale=0.025]{assets/d3.png} $\}$, $\ $ \small{Chair} $=$ \includegraphics[scale=0.02]{assets/committee.png} $\Rightarrow$ \small{Positional Scoring Rule}
		\vspace{2em}
		\begin{columns}
			\begin{column}{0.45\textwidth}
				\begin{figure}
					\begin{tabular}{ccc}
						\visible<2->{\subfloat{\includegraphics[scale=0.12]{assets/v1.png}}} &
						\visible<2->{\subfloat{\includegraphics[scale=0.12]{assets/v2.png}}} &
						\visible<2->{\subfloat{\includegraphics[scale=0.12]{assets/v3.png}}} \\
						\visible<3->{\subfloat{\includegraphics[scale=0.05]{assets/d1.png}}} &
						\visible<4->{\subfloat{\includegraphics[scale=0.05]{assets/d1.png}}} & \\
						\visible<3->{\subfloat{\includegraphics[scale=0.05]{assets/d2.png}}} &
						\visible<4->{\subfloat{\includegraphics[scale=0.05]{assets/d2.png}}} & \\
						\visible<3->{\subfloat{\includegraphics[scale=0.05]{assets/d3.png}}} & & \\
					\end{tabular}
				\end{figure}
			\end{column}
			\begin{column}{0.45\textwidth}
				\begin{center}
					\visible<5->{\includegraphics[width=0.5\textwidth]{assets/committee.png}}\\
					\vspace{2em}
					\visible<6->{weight$(1$\textsuperscript{st}$)$ $\geq$ 2 $\cdot$ weight$(2$\textsuperscript{nd}$)$}
				\end{center}
				
			\end{column}
		\end{columns}
	
\end{frame}

\begin{frame}
	\frametitle{Introducing the problem}
	\textbf{Setting}: Incompletely specified preferences and social choice rule
	\onslide<2>{
		\begin{itemize}
			\item Agents: difficult or costly to order \emph{all} alternatives
			\item Chair: difficult to \emph{specify} a voting rule precisely
		\end{itemize}
	}
	\onslide<3-> \textbf{Goal}: Reduce uncertainty, inferring (\textit{eliciting}) incrementally and simultaneously the true preferences of agents and chair to quickly converge to an optimal or a near-optimal alternative \\
	\vspace{3em}
	\onslide<4-> \textbf{Approach}:
	\begin{itemize}
		\item Develop query strategies that interleave questions to the chair and to the agents
		\item Use \emph{Minimax regret} to measure the quality of those strategies
	\end{itemize}
\end{frame}

\subsection{Related Works}
\begin{frame}
	\frametitle{Related Works}
	\textbf{Incomplete profile}  
	\begin{itemize}
		\item and known weights: Minimax regret to produce a robust winner approximation (\textit{Lu and Boutilier 2011}, \cite{Lu2011}; \textit{Boutilier et al. 2006}, \cite{Boutilier2006})
	\end{itemize}~\\
	\textbf{Uncertain weights} 
	\begin{itemize}
		\item and complete profile: dominance relations derived to eliminate alternatives always less preferred than others (\textit{Stein et al. 1994}, \cite{Stein1994})
		\item in positional scoring rules (\textit{Viappiani 2018}, \cite{Viappiani2018})
	\end{itemize}
\end{frame}


\begin{frame}
	\frametitle{Assumptions}
	\begin{itemize}
		\item We consider \textit{Positional Scoring Rules}, which attach weights to positions according to a scoring vector $W$
		\item We assume $W$ to be \textit{convex}
		\[ W_r - W_{r+1} \geq W_{r+1}-W_{r+2}\]
		for all positions $r$, and that $W_1=1$ and $W_m=0$
	\end{itemize}	
\end{frame}


\begin{frame}
	\frametitle{Notation}	
	\begin{description}[$W=(\w_k,\ 1\leq k \leq m), \ W \in \mathcal{W}$]
		\item [$A$] alternatives, $|A|=m$
		\item [$N$] agents (\textit{voters})
		\item [$P = (\pref_j, \ j \in N ), \ P \in \mathcal{P} $] complete preferences profile 
		\item [$W=(W_r,\ 1\leq r \leq m), \ W \in \mathcal{W}$] (convex) scoring vector that the chair has in mind
	\end{description}
	\bigskip
	\onslide<2-> \begin{block}{}
		$W$ defines a Positional Scoring Rule $f_W(P)\subseteq A$ using scores \usebeamercolor*[fg]{description item}$s^{W,P}(a), \ \forall \ a \in A$
	\end{block}
	\onslide<3-> \begin{block}{}
		$P$ and $W$ exist in the minds of agents and chair but unknown to us
	\end{block}
	
\end{frame}


%%%
%\begin{frame}
%	\frametitle{Outline}
%	\tableofcontents[hideallsubsections, sectionstyle=shaded/show]
%\end{frame}
%
%\AtBeginSection{
%	\begin{frame}
%		\frametitle{Outline}
%		\tableofcontents[currentsection, hideallsubsections]
%	\end{frame}
%}

\begin{frame}
	\frametitle{Questions}
	Two types of questions:\\
	\vspace{0.5cm}
	\onslide<2-> \textbf{Questions to the agents}
	\begin{itemize}
		\item[] Comparison queries that ask a particular agent $j$ to compare two alternatives $a, b \in \allalts$
		\usebeamercolor*[fg]{description item}\[a \pref_j b \quad ?\]
	\end{itemize}
	\onslide<3->  \textbf{Questions to the chair}
	\begin{itemize}
		\item[] Queries relating the difference between the importance of consecutive ranks from $r$ to $r+2$
		\usebeamercolor*[fg]{description item} \[ W_{r} - W_{r+1} \geq \lambda (W_{r+1} - W_{r+2}) \quad ? \] 
	\end{itemize}
\end{frame}

\begin{frame}
	\frametitle{Current Knowledge}
	The answers to these questions define $C_P$ and $C_W$ that is our knowledge about P and W 
	\medskip
	\begin{itemize}
		\onslide<2->{\item $C_P \subseteq \mathcal{P}$ constraints on the profile given by the agents}
		\onslide<3->{\item $C_W \subseteq \mathcal{W}$ constraints on the voting rule given by the chair}
	\end{itemize}

	\def\Pcircle{(0,0) circle (0.6cm)}
	\def\CPcircle{(0,0) circle (0.2cm)}
	\def\Wcircle{(0,-1.6) circle (0.6cm)}
	\def\CWcircle{(0,-1.6) circle (0.2cm)}
	\def\Allcircle{(3,-0.8) circle (0.6cm)}
	\def\Currcircle{(3,-0.8) circle (0.2cm)}
	\def\wcircle{(3,-0.8) circle (0.05cm)}
	\begin{center}
		\begin{tikzpicture}
			\onslide<2->{\draw \Pcircle node[left] at (0,0.7) {\footnotesize$\mathcal{P}$};
			\filldraw[fill=gray, draw=black] \CPcircle node[above]  at (0,0.07) {\footnotesize$C_P$};}
			\onslide<3->{\draw \Wcircle node[left] at (0,-0.9) {\footnotesize$\mathcal{W}$};
			\filldraw[fill=gray, draw=black] \CWcircle node[above]  at (0,-1.5) {\footnotesize$C_W$};}
			\onslide<4->{\draw [decorate, decoration = {brace, raise=40pt, amplitude=5pt}] (0,0.7) --  (0,-2.3);}
			\onslide<5->{\draw \Allcircle node[above] at (3,-0.2) {\footnotesize$\powersetz{\allalts} {\scriptstyle \ = \ 2^{\allalts}\setminus \emptyset}$ };}
			\onslide<5->{\filldraw[fill=gray, draw=black] \Currcircle node[above]  at (3,-0.7){\footnotesize$C_{\mathscr{P}}$};}
			\onslide<6->{\filldraw[fill=black, draw=black] \wcircle;}
		\end{tikzpicture}
	\end{center}
\end{frame}

\begin{frame}
	\frametitle{Minimax Regret}
	Given $C_P \subseteq \mathcal{P}$ and $C_W \subseteq \mathcal{W}$:
	
	\begin{block}{}
		\[\color{darkred}{\PMR^{C_P,C_W}(a,b)}= \max_{P\in C_P, W \in C_W} s^{P,W}(b)-s^{P,W}(a) \]
		is the maximum difference of score between $a$ and $b$ under all possible realizations of the full profile {\em and} weights
	\end{block}
	
	\onslide<2->  We care about the worst case loss: \emph{maximal regret} between a chosen alternative $a$ and best real alternative $b$
	\[\MR^{C_P,C_W}(a)= \max_{b\in \allalts} \PMR^{C_P,C_W}(a,b) \]
	 	
	\onslide<3-> \centerline{\textbf{We select the alternative which \emph{minimizes} the maximal regret}}
	\[\MMR^{C_P,C_W}= \min_{a\in \allalts} \MR^{C_P,C_W}(a)\]
\end{frame}

\begin{frame}
	\frametitle{Pairwise Max Regret Computation}
	The computation of $\PMR^{C_P,C_W}( \icimg{assets/d1.png},\icimg{assets/d3.png})$ can be seen as a game in which an adversary both:\begin{itemize}
		\onslide<2-> \item \textbf{chooses a complete profile $\mathbf{P \in \mathcal{P}}$}\\
		\begin{columns}
			\begin{column}{0.4\textwidth}
				\begin{figure}
					\begin{tabular}{ccc}
						\subfloat{\includegraphics[scale=0.12]{assets/v1.png}} &
						\subfloat{\includegraphics[scale=0.12]{assets/v2.png}} &
						\subfloat{\includegraphics[scale=0.12]{assets/v3.png}} \\
						\subfloat{\includegraphics[scale=0.05]{assets/d1.png}} &
						\subfloat{\includegraphics[scale=0.05]{assets/d1.png}} & \\
						\subfloat{\includegraphics[scale=0.05]{assets/d2.png}} &
						\subfloat{\includegraphics[scale=0.05]{assets/d2.png}} & \\
						\subfloat{\includegraphics[scale=0.05]{assets/d3.png}} & & \\
					\end{tabular}
				\end{figure}
			\end{column}
			\begin{column}{0.08\textwidth}
				\icarr{arrow.png}
			\end{column}
			\begin{column}{0.4\textwidth}
				\begin{figure}
					\begin{tabular}{ccc}
						\subfloat{\includegraphics[scale=0.12]{assets/v1.png}} &
						\subfloat{\includegraphics[scale=0.12]{assets/v2.png}} &
						\subfloat{\includegraphics[scale=0.12]{assets/v3.png}} \\
						\subfloat{\includegraphics[scale=0.05]{assets/d1.png}} &
						\subfloat{\includegraphics[scale=0.05]{assets/d3.png}} & 
						\subfloat{\includegraphics[scale=0.05]{assets/d3.png}} \\
						\subfloat{\includegraphics[scale=0.05]{assets/d2.png}} &
						\subfloat{\includegraphics[scale=0.05]{assets/d1.png}} & 
						\subfloat{\includegraphics[scale=0.05]{assets/d2.png}} \\
						\subfloat{\includegraphics[scale=0.05]{assets/d3.png}} &
						\subfloat{\includegraphics[scale=0.05]{assets/d2.png}} & 
						\subfloat{\includegraphics[scale=0.05]{assets/d1.png}} \\
					\end{tabular}
				\end{figure}
				
			\end{column}
		\end{columns}
		
		\onslide<3-> \item \textbf{chooses a feasible weight vector $\mathbf{W \in \mathcal{W}}$}\\
		\medskip
		\centerline{\color{darkred}$\mathbf{(1,?,0)} \qquad$ \icarr{arrow.png} \color{darkred}$\qquad \mathbf{(1,0,0)}$}
	\end{itemize}
	\medskip
	in order to maximize the difference of scores
\end{frame}
	
	\begin{frame}
		\frametitle{Elicitation strategies}
		At each step, the strategy selects a question to ask either to one of the agents about her preferences or to the chair about the voting rule \\ \bigskip
		\onslide<2-> The termination condition could be:
		\begin{itemize}
			\item <3-> when the minimax regret is lower than a threshold
			\item <4-> when the minimax regret is zero
		\end{itemize}
		\bigskip
	\end{frame}
	
	\begin{frame}[t]
		\frametitle{Elicitation strategies}
		\framesubtitle{Pessimistic Strategy}
		\onslide<1-> Assume that a question leads to the possible new knowledge states $(C_P^1, C_W^1)$ and $(C_P^2, C_W^2)$ depending on the answer, then the badness of the question in the worst case is:
		\[\max_{i=1,2} \MMR(C_P^i, C_W^i) \]
		\onslide<2-> The pessimistic strategy selects the question that leads to minimal regret in the worst case from a set of $n+1$ candidate questions \\
		\bigskip
		{\small \onslide<3-> \begin{block}{Note:}
			if the maximal MMR of two questions are equal, then prefers the one with the lowest MMR values associated to the opposite answer
		\end{block}}
	\end{frame}
	
	\begin{frame}
	\frametitle{Elicitation strategies}
	\framesubtitle{Pessimistic Strategy: Candidate questions}
	Let $(a^{*}, \bar{b}, \bar{P}, \bar{W})$ be the current solution of the minimax regret \\ \vspace{0.5em}
	We select $n + 1$	candidate questions:
	\begin{itemize}
		\item \textbf{One question per agent:} For each agent $i$, either:  
			\begin{itemize}
				\item $a^* \pref^{\bar{P}}_j \bar{b}$ : we ask about an incomparable alternative that can be placed above $a^*$ by the adversary to increase PMR($a^*$,$\bar{b}$)
				\item $\bar{b} \pref^{\bar{P}}_j a^*$: we ask about an incomparable alternative that can be placed between $a^*$ and $\bar{b}$ by the adversary to increase PMR($a^*$,$\bar{b}$) 
				\item $a^*$ and $\bar{b}$ are incomparable: we ask to compare them
		\end{itemize}
		\item \textbf{One question to the chair:} Consider $W_\tau$ the weight vector that minimize the PMR in the worst case. \\ We ask about the position
		$r = \argmax\limits_{i= [\![1, m-1]\!] } |\bar{W}(i)-W_\tau(i)|$
		
	\end{itemize}
	\end{frame}
	
	\begin{frame}
		\frametitle{Empirical Evaluation}
		\framesubtitle{Pessimistic for different datasets}
		\begin{figure}
			\centering
			\caption{Average MMR (normalized by $n$) after $k$ questions with Pessimistic strategy for different datasets.}
			\label{fig:linearity}
			\begin{tikzpicture}
				\pgfplotsset{
					every axis legend/.append style={
						at={(0.5,1.1)},
						anchor=south
					},
				}
				\begin{axis}[
					y=80,
					legend columns=3,
					xlabel=Number of Questions,
					ylabel=MMR/n,
					ytick={0,0.5,1},
					xtick distance=100,
					xtick pos=left,
					ymajorgrids=true,
					enlarge x limits=-1, %hack to plot on the full x-axis scale
					width=10cm, %set bigger width
					ytick style={draw=none},
					ymin=0,
					ymax=1,
					xmin=0,
					xmax=1000,
					yticklabels={0,0.5,1},
					legend style={font=\footnotesize}]
					
					\addlegendimage{mark=halfsquare right*,brown,mark size=2}
					\addlegendimage{mark=diamond*,red,mark size=2}
					\addlegendimage{mark=pentagon*,cyan,mark size=2}
					\addlegendimage{mark=halfcircle*,violet,mark size=2}
					\addlegendimage{mark=*,pink,mark size=2}
					\addlegendimage{mark=triangle*,green,mark size=2}
					\addlegendimage{mark=halfsquare left*,blue,mark size=2}
					\addlegendimage{mark=square*,teal,mark size=2}
					\addlegendimage{mark=halfsquare*,magenta,mark size=2}
					
					
					\addplot[thick, mark=halfsquare right*, mark size = {2}, mark indices = {120}, brown] table [x=k, y=5.20]{data/linearity.dat};
					\addlegendentry{m=5, n=20}
					\addplot[thick, mark=diamond*, mark size = {2}, mark indices = {150}, red] table [x=k, y=10.20]{data/linearity.dat};
					\addlegendentry{m=10, n=20}
					\addplot[thick, mark=pentagon*, mark size = {2}, mark indices = {240}, cyan] table [x=k, y=11.30]{data/linearity.dat};
					\addlegendentry{m=11, n=30}
					\addplot[thick, mark=halfcircle*, mark size = {2}, mark indices = {400}, violet] table [x=k, y=tshirts]{data/linearity.dat};
					\addlegendentry{tshirts m11n30}
					\addplot[thick, mark=*, mark size = {2}, mark indices = {400}, pink] table [x=k, y=courses]{data/linearity.dat};
					\addlegendentry{courses m9n146}
					\addplot[thick, mark=triangle*, mark size = {2}, mark indices = {400}, green] table [x=k, y=9.146]{data/linearity.dat};
					\addlegendentry{m=9, n=146}
					\addplot[thick, mark=halfsquare left*, mark size = {2}, mark indices = {200}, blue] table [x=k, y=14.9]{data/linearity.dat};
					\addlegendentry{m=14, n=9}
					\addplot[thick, mark=square*, mark size = {2}, mark indices = {60}, teal] table [x=k, y=skate]{data/linearity.dat};
					\addlegendentry{skate m14n9}
					\addplot[thick, mark=halfsquare*, mark size = {2}, mark indices = {400}, magenta] table [x=k, y=15.30]{data/linearity.dat};
					\addlegendentry{m=15, n=30}
				\end{axis}
			\end{tikzpicture}
		\end{figure}
	\end{frame}
	
	\begin{frame}
		\frametitle{Empirical Evaluation}
		\framesubtitle{Pessimistic reaching "low enough" regret}
		\sisetup{table-number-alignment = center, table-figures-integer=2, table-figures-decimal=1, table-auto-round}
		\begin{minipage}{\textwidth}
			\captionof{table}{Questions asked by Pessimistic strategy on several datasets to reach $\frac{n}{10}$ regret, columns 4 and 5, and zero regret, last two columns.}
			\label{tab:questions}
			\begin{center}
			\scalebox{0.8}{
			\begin{tabular}{cccS[table-number-alignment = center, table-figures-integer=2] S[table-figures-integer=3, table-figures-decimal=1]@{ $\vert$ }S[table-figures-integer=2, table-figures-decimal=1]@{ $\vert$ }S[table-figures-integer=2, table-figures-decimal=1]@{ ]} S[table-number-alignment = center, table-figures-integer=4]S[table-figures-integer=3, table-figures-decimal=1]@{ $\vert$ }S[table-figures-integer=2, table-figures-decimal=1]@{ $\vert$ }S[table-figures-integer=2, table-figures-decimal=1]@{ ]}}
				\toprule
				{dataset} & m & n &{$q_{c}^{\scriptscriptstyle{MMR} \leq n/10}$} & \multicolumn{3}{c}{$q_{a}^{\scriptscriptstyle{MMR} \leq n/10}$} & {$q_{c}^{ \scriptscriptstyle{MMR} = 0}$} & \multicolumn{3}{c}{$q_{a}^{ \scriptscriptstyle{MMR} = 0}$} \\
				\midrule
				m5n20 & 5&20&0.0&[4.3 &4.95 & 5.84 &5.25&[ \ 5.36 & 6.15 & 7.21\\
				m10n20&10&20&0.0&[13.85 & 16.1& 18.41&31.95&[19.66 & 21.78 & 24.7\\
				m11n30&11&30&0.0&[16.55&19.0&22.26&45.15&[23.07&25.7&28.89\\
				tshirts&11&30&0.0&[13.08&16.6&19.58& 43.15 &[28.22&31.98 &35.62\\
				courses&9&146&0.0&[6.03 &7.0 &7.0&0.0&[6.81 & 7.0 &7.0\\
				%m9n146&9&146&0&0&[1.94 &8 &9.25&999.3&0.47&0&0&[1.94&8&9.25&999.3&0.47\\
				m14n9&14&9&5.4&[30.3&33.45&36.65&64.05&[37.55&40.5&44.3\\
				skate&14&9&0.0&[11.35&11.6&12.3&0.0&[11.5&11.8&12.75 \\
				m15n30&15&30&0.0&[24.95&29.5&33.68 \\
				\bottomrule
			\end{tabular}}
		\end{center}
		\end{minipage}
	\end{frame}


\begin{frame}
	\frametitle{Empirical Evaluation}
	\framesubtitle{Pessimistic chair first and then agents (and vice-versa)}
	\sisetup{table-number-alignment = center, table-figures-integer=2, table-figures-decimal=1, table-auto-round}
	\begin{minipage}{\textwidth}
		\centering
		\captionof{table}{Average MMR in problems of size $(10, 20)$ after $500$ questions, among which $q_c$ to the chair.}
		\label{tab:twoP500}
		\scalebox{0.85}{
			\begin{tabular}{S[table-figures-integer=3, table-figures-decimal=0]S[table-number-alignment = right]@{ ± }S[table-number-alignment = left, table-figures-integer=1]S[table-number-alignment = right]@{ ± }S[table-number-alignment = left, table-figures-integer=1]}
				\toprule
				{$q_c$} & {2 ph.\ ca} & {sd} & {2 ph.\ ac} & {sd} \\
				\midrule		
				
				0	&	0.62	&	0.52	&	0.62	&	0.52	\\
				15	&	0.515	&	0.48	&	0.54	&	0.46	\\
				30	&	0.345	&	0.47	&	0.325	&	0.425	\\
				50	&	0.045	&	0.09	&	0.03	&	0.065	\\
				100	&	0.14	&	0.23	&	0.075	&	0.135	\\
				200	&	2.305	&	1.36	&	2.145	&	1.845	\\
				300	&	5.15	&	2.38	&	6.83	&	0.625	\\
				400	&	10.905	&	0.89	&	12.245	&	0.99	\\
				500	&	20.0	&	0.0	&	20.0	&	0.0	\\
				
				\bottomrule
			\end{tabular}
		}
	\end{minipage}
\end{frame}



\addtocounter{framenumber}{-1}
\begin{frame}[plain]
	\centering {\color{darkred}\LARGE Thank You!}
	{\let\thefootnote\relax\footnote{\tiny Icon made by Eucalyp from www.flaticon.com}}
\end{frame}

\bibliographystyle{plain}
\begin{frame}[plain,allowframebreaks,noframenumbering]
	\bibliography{biblio}
\end{frame}

\begin{frame}[plain,noframenumbering]
	\frametitle{Building concrete questions for the chair}
	Queries relating the difference between the importance of consecutive ranks
	\only<1>{{\usebeamercolor*[fg]{description item} \[ W_{r} - W_{r+1} \geq \lambda (W_{r+1} - W_{r+2}) \quad ? \] }}
	\onslide<2->{{\usebeamercolor*[fg]{description item} \[ W_{2} - W_{3} \geq 2 \ (W_{3} - W_{4}) \quad ? \] }}
	\onslide<3->{Abstract queries are difficult to answer}
	{\usebeamercolor*[fg]{description item}
	\begin{align*}
		\onslide<4->{W_{2}\ - \ W_{3} &\geq 2 \ W_{3} - 2 \ W_{4} }\\
		\onslide<5->{W_{2} + 2 \ W_{4} &\geq 3 \ W_{3} }\\
		\onslide<6->{s(a) &\geq s(b)}
	\end{align*}
	}
	\only<7>{\begin{figure}
		\begin{center}
			$
			\begin{array}{cc}
				\mathbf{\#1}&\mathbf{\#2} \\
				-&-\\
				\mathbf{a}&-\\
				\mathbf{b}&\mathbf{b}\\
				-&\mathbf{a}\\
			\end{array}
			$
		\end{center}
	\end{figure}
	}
	\onslide<8->{
	\begin{figure}
		\begin{center}
			$
			\begin{array}{cccc}
				\mathbf{\#1}&\mathbf{\#2}&\mathbf{\#3}&\mathbf{\#3} \\
				c&d&a&b\\
				\mathbf{a}&c&b&a\\
				\mathbf{b}&\mathbf{b}&c&d\\
				d&\mathbf{a}&d&c\\
			\end{array}
			$
		\end{center}
	\end{figure}
	}
\end{frame}


\end{document}